\documentclass[11pt]{article}

\usepackage[utf8]{inputenc}
\usepackage[T1]{fontenc}
\usepackage{microtype}
\usepackage[margin=1in]{geometry}
\usepackage{graphicx}
\usepackage{booktabs}
\usepackage{amsmath}
\usepackage{amssymb}
\usepackage{xcolor}
\usepackage[numbers,sort&compress]{natbib}
\usepackage{hyperref}
\hypersetup{colorlinks=true, linkcolor=blue!55!black, citecolor=blue!55!black,
            urlcolor=blue!55!black}
\usepackage{authblk}
\usepackage{caption}
\captionsetup{font=small, labelfont=bf}

\title{Deciphering the molecular heterogeneity and dynamics of human
hippocampal neural stem cells across neonatal, adult, aging, and
stroke-injured states}

\author[]{}
\affil[]{}
\date{}

\begin{document}

\maketitle

\begin{abstract}
Whether the adult human hippocampus retains the capacity for neurogenesis has
been debated for decades, and the molecular states of its neural stem cells
(NSCs) across life stages and after injury remain unclear. Here we generate a
single-nucleus RNA sequencing atlas of the anterior-mid hippocampus from
ten postmortem donors spanning neonatal (postnatal day 4, $n=1$), adult
(31 and 32 years, $n=2$), aging (50--68 years, $n=6$), and stroke-induced
injury (48 years, $n=1$) contexts. After quality control we retained 92{,}966
of 99{,}635 sequenced nuclei and resolved 16 major populations by Uniform
Manifold Approximation and Projection (UMAP), including astrocyte1 (1{,}146),
astrocyte2/quiescent NSC (AS2/qNSC, 11{,}071), primed NSC (pNSC, 2{,}798),
active NSC (aNSC, 2{,}140), neuroblast (2{,}607), and granule cell (24{,}671).
Subclustering the AS2/qNSC compartment separated astrocytes from deep
quiescent qNSC1 (HOPX-high) and shallow quiescent qNSC2 (LPAR1-high).
Combining single-cell hierarchical Poisson factorization with Seurat marker
analysis yielded human-specific lineage markers --- LRRC3B, RHOJ, SLC4A4 and
GLI3 for quiescent NSCs; CHI3L1 and EGFR for pNSCs; and CALM3, TTC9B, NRGN,
FXYD7, NRN1, GNG3, TCEAL5, TMSB10 and NEUROD2 for neuroblasts --- that avoid
the GABAergic-interneuron contamination confounding DCX, PROX1 and STMN2.
RNA-velocity and Monocle pseudotime analyses in the neonatal hippocampus
established a pNSC $\rightarrow$ aNSC $\rightarrow$ neuroblast trajectory that
bifurcates into two granule-cell subtypes, validated by HOPX, NES, Ki67 and
PSA-NCAM immunostaining. With aging, most NSCs entered deep quiescence and
pNSC/aNSC populations were almost absent. After stroke, by contrast, pNSCs and
aNSCs reached 17.3\% (322/1861) and 11.7\% (218/1861) of the neurogenic lineage
--- comparable to the neonatal cohort and well above adult and aged samples ---
consistent with injury-induced reactivation of dormant NSCs. Together, these
results unify conflicting reports of human adult hippocampal neurogenesis by
showing that a reserved quiescent pool persists into old age and can be
re-awakened by pathological signals.
\end{abstract}

\section{Introduction}

Lifelong learning and memory consolidation rest in part on continued
neurogenesis in the dentate gyrus (DG) of the hippocampus, where resident
neural stem cells (NSCs) generate granule cells throughout development and,
in rodents, into adulthood. The developmental logic worked out in mice places
radial glia-like (RGL) cells at the apex of a lineage that passes through
quiescent and active intermediate states, through neuroblasts, and on to
granule-cell neurons, with each transition governed by conserved signalling
programs \cite{encinas2011division,obernier2019neural,urban2019quiescence}.
How, and whether, this machinery operates in the adult \emph{human}
hippocampus has been contested. Immunohistochemical studies have reached
opposite conclusions: some describe a sharp collapse of DCX-positive cells
after the first year of life, with essentially no young neurons detectable in
adults \cite{sorrells2018human,cipriani2018hippocampal,franjic2022transcriptomic},
whereas others document thousands of immature neurons into the eighth and
ninth decade, sometimes increased in the vicinity of neurodegenerative
pathology \cite{boldrini2018human,morenojimenez2019adult,tobin2019human,terrerosroncal2021impact}.
Each side marshals plausible technical arguments, and neither has been able to
resolve the question with the cell-type specificity that single-cell
transcriptomics affords.

Single-nucleus RNA sequencing (snRNA-seq) has become the method of choice for
human brain tissue because it operates on frozen postmortem material, avoids
antibody-specificity pitfalls, and resolves populations that are rare enough
to be invisible to bulk assays \cite{habib2017massively}. Recent atlases have
sharpened our view of the developing \cite{zhong2020decoding} and adult
\cite{franjic2022transcriptomic,ayhan2021resolving} hippocampus, yet they
leave the molecular physiology of NSCs themselves --- how the quiescent,
primed, and active states evolve across life stages, and whether insult can
bring dormant cells back into the cycle --- largely unresolved. In the
rodent subventricular zone, ischemia reactivates a dormant NSC subpool
through interferon-$\gamma$--linked priming \cite{llorensbobadilla2015single};
the equivalent question in the human hippocampus has not been addressed with
transcriptomic resolution, in part because donor brains covering injury
states are exceptionally rare.

We set out to fill this gap with an snRNA-seq atlas of the anterior-mid
hippocampus from ten postmortem donors that span four biologically distinct
contexts: neonatal (postnatal day 4, $n=1$), adult (31- and 32-year-old,
$n=2$), aging (50, 56, 60, 64, 64, and 68-year-old, $n=6$), and stroke-induced
injury (48-year-old, $n=1$). The donor design deliberately places a pediatric
baseline, two healthy adult windows and a single acute-injury sample alongside
each other so that within-lineage composition changes and transcriptomic
responses can be tracked on a shared coordinate system. We then combine
cross-species integration with mouse, pig and macaque data
\cite{hochgerner2018conserved,franjic2022transcriptomic}, single-cell
hierarchical Poisson factorization \cite{levitin2019gene}, Seurat-based
marker discovery and integration \cite{butler2018integrating}, RNA-velocity
\cite{lamanno2018rna}, and Monocle~2 pseudotime reconstruction
\cite{trapnell2014dynamics}, with paired immunostaining, to unify several
questions that have until now been approached separately.

Our analysis yields four principal findings. First, the AS2/qNSC compartment
in the neonatal hippocampus cleanly divides into astrocytes, deep-quiescent
qNSC1, and shallow-quiescent qNSC2, each with distinct transcriptomic and
cell-cycle properties. Second, single-cell hierarchical Poisson factorization
plus Seurat marker discovery, followed by explicit filtering against the
GABAergic-interneuron transcriptome, produces a panel of human-specific
neurogenic-lineage markers (LRRC3B, RHOJ, SLC4A4 and GLI3 for qNSCs; CHI3L1
and EGFR for pNSCs; and CALM3, TTC9B, NRGN, FXYD7, NRN1, GNG3, TCEAL5,
TMSB10 and NEUROD2 for neuroblasts) that addresses the non-specificity of
DCX, PROX1 and STMN2 in human tissue. Third, pseudotime reconstruction in
the neonatal DG recovers a pNSC $\rightarrow$ aNSC $\rightarrow$ neuroblast
trajectory that bifurcates into two granule-cell end-states, with transcription
factors and marker cascades that agree with paired HOPX, NES, Ki67 and
PSA-NCAM immunostaining. Fourth, compared to adult and aged hippocampi in
which pNSC/aNSC populations are essentially absent, the stroke-injured
donor shows pNSC and aNSC fractions of 17.3\% (322/1861) and 11.7\%
(218/1861) of the neurogenic lineage --- comparable to the neonatal cohort
--- supporting the interpretation that injury can reactivate dormant
quiescent NSCs in humans. Taken together, these results provide a concrete
molecular basis for the recurring observation that adult human hippocampal
neurogenesis is difficult to detect under baseline conditions but can be
caught in the act under disease or injury.

\section{Methods}

\subsection{Donor cohort and tissue collection}

We recruited ten postmortem donors spanning four life-stage and health
contexts: neonatal (postnatal 4 days, $n=1$); adult (31 and 32 years, $n=2$);
aging (50, 56, 60, 64, 64 and 68 years, $n=6$); and stroke-induced injury
(48 years, $n=1$). Nine donors were male and one was female. For each hippocampus
we collected anterior and mid segments containing a typical dentate gyrus
structure, as the rostro-caudal axis of the primate hippocampus shows
heterogeneous cell composition that would otherwise confound across-donor
comparison. A cynomolgus monkey (3 months, 2.3 kg, female) served as a
primate comparator for immunostaining. Tissue was snap-frozen in OCT
(Tissue-Tek) on dry ice and sectioned at 10 $\mu$m on a Leica CM1950
cryostat for staining, with a paired portion retained for nuclei isolation.

\subsection{Nuclei isolation and snRNA-seq}

We followed an snRNA-seq-compatible protocol similar to that described for
human hippocampal tissue \cite{ayhan2021protocol,habib2017massively}. Frozen
tissue was placed in chilled lysis buffer containing 10~mM Tris-HCl, 10~mM
NaCl, 3~mM MgCl$_2$ and 0.1\% Nonidet P40 Substitute (Sigma Aldrich,
PN-74385) and incubated on ice for 15 minutes with gentle shaking. The
homogenate was passed through a 30~$\mu$m MACS SmartStrainer (Miltenyi
Biotec, PN-130-098-458) to remove debris and aggregates, centrifuged at
500$\times$g for 5 minutes at 4$^\circ$C, and resuspended in Nuclei Wash and
Resuspension Buffer (1$\times$ PBS, 1.0\% BSA, 0.2~U/$\mu$L RNase inhibitor;
Sigma Aldrich PN-3335399001), then gently pipetted eight to ten times.
Nuclei were loaded onto a 10x Genomics Chromium controller and prepared
libraries were sequenced to yield 99{,}635 raw nuclei.

\subsection{Quality control and clustering}

We removed cell debris, cell aggregates, and cells with more than 20\% of
transcripts arising from mitochondrial genes, retaining 92{,}966 nuclei at a
median of 3{,}001 genes per nucleus. Clustering and visualization used
Uniform Manifold Approximation and Projection (UMAP) within the Seurat
framework \cite{butler2018integrating}. Marker genes were called with the
Wilcoxon rank-sum test (Seurat default), selecting genes with average log-expression
differences greater than 0.5 and p-values less than 0.05. Gene ontology
enrichment of Biological Process terms used the clusterProfiler package
with p $<$ 0.05 for each cell population.

\subsection{Cell-type specific analyses}

To distinguish quiescent NSCs from astrocytes, which share GFAP, ALDH1L1 and
VIM, we subclustered the AS2/qNSC population with Seurat FindAllMarkers and
annotated subclusters by gene set scores (AddModuleScore) built from published
astrocyte and RGL panels \cite{zamanian2012genomic,liddelow2017neurotoxic,clarke2018normal,hochgerner2018conserved,zhong2020decoding,franjic2022transcriptomic}.
Cell-cycle state was estimated with Seurat's CellCycleScoring.
To mine novel markers for human neurogenic populations we combined
single-cell hierarchical Poisson factorization (scHPF)
\cite{levitin2019gene}, which identifies specific rather than merely abundant
genes by modelling count sparsity, with Seurat FindAllMarkers for highly
expressed genes. Candidate neuroblast markers were then mapped onto the
GABAergic-interneuron population, and genes with low or absent expression
in interneurons were retained.

\subsection{Lineage trajectory reconstruction}

RNA-velocity analysis \cite{lamanno2018rna} was applied to the neonatal and
injury-sample neurogenic lineages to estimate directional transitions.
Pseudotime reconstruction used Monocle~2 \cite{trapnell2014dynamics}: dispersed
genes were selected with the \texttt{estimateDispersions} function and the
\texttt{DDRTree} method was used for dimensionality reduction
(\texttt{reduceDimension(mycds, max\_components = 2, method = "DDRTree")}),
with the minimum spanning tree visualized via \texttt{plot\_cell\_trajectory}.
Branched expression analysis modelling (BEAM) was used to recover genes that
differentiate between branches in the neonatal sample, retaining genes at a
q-value threshold below $1\times 10^{-4}$, which yielded four gene clusters.

\subsection{Immunostaining and TUNEL}

Hippocampal tissue was fixed in 4\% paraformaldehyde for up to 24 hours,
cryoprotected in 30\% sucrose at 4$^\circ$C, and sectioned at 10 $\mu$m.
Sections were blocked and permeabilized in 2\% Triton X-100 for 2 hours,
treated with a VECTOR TrueView autofluorescence quenching kit
(Vectorlabs, PN-SP-8500-15), washed in PBS (3$\times$15 minutes), and
incubated in 3\% BSA for one hour at room temperature. After overnight
primary-antibody incubation, sections were washed (PBS, 3$\times$15 minutes)
and incubated with Alexa Fluor 488 or Alexa Fluor 647 AffiniPure donkey
anti-rabbit IgG(H+L) (1:500; Jackson ImmunoResearch, PN-712-545-152 or
PN-715-605-150) for 2 hours at room temperature, with Alexa Fluor 555
donkey anti-rabbit IgG (1:500; ThermoFisher PN-A-31572) used where
appropriate, and DAPI (Sigma, PN-32670-5mg-F) for nuclear counterstain.
To label damaged nuclei, 20~$\mu$L of TUNEL reaction mixture (labelling
buffer, TdT, BIO-d-UTP) was applied for 120 minutes at 37$^\circ$C in a
humidified chamber after a 2-minute PBS wash and 30-minute blocking in
50~$\mu$L of blocking reagent. Apical-process-bearing radial-glia
morphology was required for scoring a cell as an active NSC.

\section{Results}

\subsection{A single-nucleus atlas of the human hippocampus across ages and
injury}

We first asked whether an snRNA-seq atlas assembled across neonatal, adult,
aging, and stroke-injured donors would recover a reproducible lineage of
neurogenic cells. To generate the atlas, we dissected anterior-mid
hippocampus (which contains a clear dentate-gyrus structure) from ten
postmortem donors (Figure~\ref{fig:cohort}) and processed the material for
10x Genomics snRNA-seq in parallel with paired immunostaining on the
contralateral tissue. Of 99{,}635 sequenced nuclei, 92{,}966 were retained
after removal of debris, aggregates, and cells with more than 20\%
mitochondrial-gene transcripts, at a median of 3{,}001 genes per nucleus.
Pooling across donors and projecting via UMAP yielded 16 major cell
populations (Figure~\ref{fig:populations}): astrocyte1 (1{,}146),
AS2/qNSC (11{,}071), pNSC (2{,}798), aNSC (2{,}140), neuroblast (2{,}607),
granule cell (24{,}671), GABAergic interneuron (8{,}601), pyramidal neuron
(676), oligodendrocyte progenitor (5{,}396), oligodendrocyte (15{,}796),
microglia (11{,}823), endothelial cell (1{,}232), pericyte (981), Reelin-
expressing Cajal--Retzius cell (218), and two unidentified populations
(UN1 and UN2 combined, 3{,}810). Each population was defined by a
combination of classical markers and differentially expressed genes, and
the average number of detected genes per cell was comparable across
life-stage groups, which rules out a gene-count artefact driving apparent
enrichment of stem-cell transcripts in specific populations.

Quantifying population fractions across life stage showed substantial
composition changes. pNSC and aNSC counts fell sharply with aging but
partially recovered in the stroke-injured hippocampus, while granule cells
were depleted in the injured donor consistent with ischemic cell loss.
These composition shifts --- and the retention of cell-cycle and
stem-cell-gene distinctions across groups --- motivated the NSC-focused
analyses that follow.

\subsection{Heterogeneity of human hippocampal NSCs}

Having established the atlas, we next asked whether the NSC compartment
resolves into transcriptomically distinct states in humans. Cross-species
integration with mouse, pig, and rhesus-macaque snRNA-seq data placed human
AS2/qNSC and pNSC with radial glia-like cells and astrocytes of the other
species, and positioned human aNSC and neuroblast with the corresponding
neural intermediate progenitors and neuroblasts, confirming the lineage
identities. Because initial UMAP clustering could not distinguish qNSCs
from astrocytes --- both express GFAP, ALDH1L1, and VIM
\cite{zamanian2012genomic,liddelow2017neurotoxic} --- we subclustered the
AS2/qNSC compartment with Seurat FindAllMarkers and obtained three groups:
astrocyte2, qNSC1, and qNSC2. Gene-set-score analysis using published
astrocyte and RGL panels \cite{clarke2018normal,hochgerner2018conserved}
showed that astrocyte2 scored highest on the astrocyte panel, whereas the
two qNSC subclusters did not; this corroborates the marker-gene separation
and places qNSC1 and qNSC2 clearly within the neurogenic lineage.

The two qNSC subclusters were not equivalent. qNSC1 preferentially expressed
the quiescence marker HOPX and lacked the neurogenic gene panel
characteristic of qNSC2, which expressed LPAR1 and neurogenic genes STMN1,
PROX1, SIRT2 and ST18, suggesting shallow quiescence with latent
lineage-progression potential. Compared with both astrocytes and qNSCs,
pNSCs showed low ALDH1L1 and GFAP but high stem-cell genes HOPX, VIM,
SOX2, SOX4 and CCND2; GO-term analysis of the top 1{,}000 pNSC-expressed
genes returned stem-cell differentiation, Wnt, Notch, Hippo, and
neurogenesis programs. aNSCs shared the stem-cell genes with pNSCs but
additionally expressed SOX11 and were enriched for neuron differentiation,
axonogenesis, and synapse-organization programs. Cell-cycle scoring
corroborated these distinctions: qNSC1 and qNSC2 were mainly in G0/G1,
whereas aNSCs concentrated in the S/G2/M phases of active mitosis. These
results together establish that the apparent AS2/qNSC mass is a compound
cluster composed of deeply quiescent, shallowly quiescent, and astrocytic
cells that can be separated transcriptomically, and that pNSCs and aNSCs
represent a graded activation axis within the primed/active compartment.

\subsection{Novel human-specific markers for NSCs and neuroblasts}

Building on this classification, we sought lineage markers that function
specifically in human tissue. To avoid the non-specificity that has dogged
DCX, PROX1, CALB2 and STMN2 in human DG, we combined single-cell
hierarchical Poisson factorization \cite{levitin2019gene} with Seurat
FindAllMarkers: scHPF isolates genes that are specifically, rather than
merely highly, expressed by a population, and FindAllMarkers
complementarily surfaces high-abundance markers. The intersection produced
a panel in which LRRC3B, RHOJ, SLC4A4 and GLI3 were specific for qNSC1 and
qNSC2, CHI3L1 and EGFR marked pNSCs, and NRGN, NRN1, and HECW1 marked
neuroblasts at the transcript level. Feature plots confirmed EGFR
specificity to pNSCs and NRGN/NRN1 specificity to neuroblasts, but also
revealed that HECW1 and several additional candidates (STMN2, NSG2, SNAP25,
BASP1) were widely expressed in the human GABAergic-interneuron population,
validated by canonical interneuron markers (DLX1, GABRG3, CCK, SLC6A11,
SLC6A1, GAD1, GAD2, CNR1, GRM1, RELN, VIP).

The overlap between neuroblast-candidate genes and interneuron expression
mirrors a recent observation for DCX \cite{franjic2022transcriptomic} and
was replicated across datasets: our human GABAergic interneurons reached
cross-species similarity scores above 0.95 to well-annotated interneurons
in mouse and macaque datasets. Canonical immature-granule-cell markers
reported in other studies, including DCX, PROX1, and CALB2 --- as well as
our own scHPF/FindAllMarkers candidates HECW1, STMN2, NSG2, SNAP25 and
BASP1 --- were therefore also enriched in the neonatal GABAergic
interneurons (and prominently so in adult interneurons). We confirmed this
at the protein level by co-immunostaining DCX with the interneuron marker
SST, which showed that somatostatin-positive interneurons in primate
hippocampus can be labelled by DCX, independent of their neurogenic
status. These findings imply that the identification of newborn neurons
in the adult human hippocampus must explicitly exclude interneuron
contamination, and motivated us to filter our neuroblast candidates
against the interneuron transcriptome.

Projecting the top-100 scHPF genes and the FindAllMarkers neuroblast set
(p-adjusted $<$ 0.01) onto the interneuron population, we retained those
genes that clustered near the expression origin in interneurons --- i.e.,
low or absent interneuronal expression --- as human neuroblast-specific
markers. This yielded CALM3, TTC9B, NRGN, FXYD7, NRN1, GNG3, TCEAL5,
TMSB10 and NEUROD2. All nine markers replicated their specificity in
adult and aging samples. In short, our combined scHPF $+$ FindAllMarkers
pipeline, coupled with an explicit interneuron-contamination filter,
produces a human-specific marker panel that can be deployed in future
studies of adult hippocampal neurogenesis.

\subsection{Developmental trajectory of NSCs in the neonatal hippocampus}

To chart developmental potential at the beginning of life, we applied
RNA-velocity \cite{lamanno2018rna} to the day-4 hippocampus. The velocity
field showed that pNSCs give two developmental directions: one into the
quiescence/astrocyte2 basin, and the other into aNSCs, consistent with the
GO signatures of the pNSC and qNSC subtypes reported above. We set pNSC as
the Monocle~2 \cite{trapnell2014dynamics} pseudotime root, because qNSCs
exit the cell cycle and development, and recovered a trajectory that
proceeds from pNSC through aNSC through neuroblast and branches into two
granule-cell neuron end-states N1 and N2. N1 cells were distinguished by
NCKAP5, SGCZ, DCC, and FAM19A2, whereas N2 cells specifically expressed
FLRT2, RIMS2, NKAIN2 and XKR4, indicating that the two branches represent
distinct granule-cell subtypes.

Along the pseudotime, HOPX, SOX2 and VIM were preferentially expressed in
pNSCs and aNSCs and declined with differentiation, whereas the neuroblast
genes DCX and STMN2 rose along the trajectory and the mature-granule-cell
gene SYT1 peaked at the terminus. Paired immunostaining confirmed the
transcriptomic picture: HOPX$^+$NES$^+$ RGLs and NES$^+$Ki67$^+$
proliferating NSCs with long apical processes were detected in the
granule-cell layer and hilus of the D4 dentate gyrus, and PSA-NCAM$^+$
neuroblasts were localized across the granule-cell layer in neonatal
clusters.

Clustering DEGs along the trajectory into four groups recovered a coherent
developmental cascade. Cluster 1, at the trajectory root, contained the
pNSC/aNSC-associated genes VIM, GFAP, SOX6, GPC6, CD44, CHI3L1, TNC, EGFR
and HOPX, enriched for cell proliferation, regeneration, angiogenesis, and
canonical Wnt signalling; HOPX, VIM, CHI3L1 and TNC were all down-regulated
along pseudotime. Cluster 2 captured mid-trajectory neurogenesis and
neuronal differentiation (NRGN, STMN2, NEUROD6). Cluster 3 (branch N2)
and Cluster 4 (branch N1) matched the GO signatures of axon guidance,
neurotransmitter secretion, long-term synaptic potentiation and ion
transport (SV2B, SYT6, SYN2) versus neuron projection guidance, dendritic
spine organization and excitatory postsynaptic potential (MYT1, GIRK2),
respectively. Transcription-factor dynamics along pseudotime showed that
progenitor-regulation factors PBX3, PROX1, GLIS3, RFX4 and TEAD1
concentrated at the root, while differentiation TFs POU6F2, FOXP2, THRB,
ETV1, NR4A3, BCL11A, NCALD, LUZP2 and RARB dominated the middle and end of
the trajectory. Together, these data establish a staged pNSC $\rightarrow$
aNSC $\rightarrow$ neuroblast $\rightarrow$ granule-cell lineage in the
neonatal human hippocampus with reproducible gene-expression and
transcription-factor signatures.

\subsection{Most NSCs become deeply quiescent during aging}

With the lineage architecture established, we next quantified its behaviour
across ages. Counting NSC subtypes and their progeny across neonatal
(D4), adult (31--32 years), and aging (50--68 years) groups showed that
the qNSC1 and qNSC2 ratios rose significantly with age, particularly the
deep-quiescent qNSC1, while pNSC and aNSC populations declined sharply
from neonatal to adult and aging stages. The neuroblast count was
comparable between neonatal and adult groups but markedly reduced in the
aging hippocampus, echoing the longer neuroblast-maturation window
reported in nonhuman primates \cite{kohler2011maturation}. Although the
aged sample contained few neurogenic cells overall, those cells that
remained still expressed the expected stem-cell and neuroblast markers
(VIM, PAX6, SOX2, PROX1, NRGN, INPP5F, TERF2IP), arguing against
contamination by other populations. Immunostaining of D4, 32y, 50y and
56y dentate gyrus corroborated these trends: NSC markers HOPX, VIM, NES
and CHI3L1 were lost in aged tissue, VIM$^+$ and NES$^+$ RGLs were only
present around the granule-cell layer and hilus at D4, and PSA-NCAM$^+$
cells with neuronal morphology were seen only in D4 and 32y samples.

To probe whether NSCs deepen their quiescence along aging, we compared
qNSCs from neonatal and aged donors. Cell-proliferation and growth-inhibition
genes (BRINP1, CABLES1, TENM2, CNTN1) and stem-cell differentiation genes
(RANBP3L, NDRG2) were up-regulated in aged qNSC1, with oligodendrocyte
genes (MBP, PLP1, MOBP) additionally up-regulated in aged qNSC2.
Meanwhile, stem-cell and regeneration genes (LPAR1, TNC, CASC15, SOX2,
SOX4, HOPX, VIM) were down-regulated in both qNSC1 and qNSC2. GO-term
enrichment of up-regulated genes pointed to negative regulation of growth,
neuronal stem-cell population maintenance, astrocyte and oligodendrocyte
differentiation, aging, and amyloid precursor protein catabolic process,
whereas the down-regulated genes enriched for nervous-system development,
neurogenesis, positive regulation of the mitotic cell cycle, tissue
regeneration, autophagy, generation of neurons, and Hippo, Wnt and Notch
signalling pathways. In agreement, comparisons of the remaining pNSC,
aNSC, and NB populations across age groups returned repression of NSC
proliferation, deficient autophagy and proteasomal protein catabolic
process, and increased glial-cell differentiation. These signatures are
consistent with the `A1-like' astrocyte-reactivity program that has been
reported in aged mouse brain \cite{clarke2018normal,liddelow2017neurotoxic}
and with the division-coupled astrocyte differentiation model of NSC
depletion developed in rodents \cite{encinas2011division}. Our results
extend those models to human tissue and imply that most human NSCs enter
deep quiescence during aging rather than being eliminated outright.

\subsection{Stroke-induced reactivation of quiescent NSCs}

Finally, because the dormant-NSC activation program has been characterized
in the mouse subventricular zone after ischemia
\cite{llorensbobadilla2015single}, we asked whether the stroke-afflicted
48-year-old donor showed equivalent signatures in the human hippocampus.
Compared with adult donors, apoptosis-, DNA-damage-, and autophagy-gene
expression was significantly up-regulated in the granule cells and
GABAergic interneurons of the injured hippocampus, and TUNEL-positive
nuclei were detected in the injured dentate gyrus but not in other
adult or aging samples, confirming injury. Coincident with this damage,
pNSC and aNSC populations reappeared and qNSCs declined, suggesting that
quiescent NSCs may have been reactivated. Immunostaining recovered
NES$^+$Ki67$^+$, NES$^+$VIM$^+$ and NES$^+$CHI3L1$^+$ active NSCs with
radial-glia morphology in the injured dentate gyrus --- cells that were
absent in all other adult samples (32y, 50y, 56y). However, because
VIM$^+$CHI3L1$^+$ reactive astrocytes also populate the injured
hippocampus, we had to separate true reactivated NSCs from reactive
astrocytes. We integrated neonatal pNSCs and aNSCs with the injury
sample as prototypes, increased clustering resolution to obtain eight
subclusters, and used gene-set scoring to assign RGL- versus
reactive-astrocyte identities. Subclusters 0, 1 and 3 were neonatal-
dominated and scored high on RGL potential, subcluster 2 was
injury-dominated and scored high on reactive-astrocyte features, and
subcluster 4 was injury-dominated but intermediate. The RGL subclusters
specifically expressed VIM, HOPX, LPAR1 and SOX2, with subcluster 0
additionally expressing STMN1, DCX and SIRT2, while the reactive-astrocyte
subcluster was defined by OSMR, TIMP1 and LGALS3. Subcluster 0 was
therefore annotated as pure aNSCs, subclusters 1 and 3 as pNSCs, and
subcluster 2 as reactive astrocytes.

With reactive astrocytes excluded, the injury-induced neurogenic lineage
shows a striking rebalancing (Figure~\ref{fig:lineage}, Figure~\ref{fig:injury_counts}).
In the stroke-injured hippocampus, pNSCs and aNSCs reached 17.3\%
(322/1861) and 11.7\% (218/1861) of the neurogenic lineage, comparable to
the neonatal cohort and well above the adult and aging groups.
Correspondingly, qNSC1 and qNSC2 fractions dropped to 23.4\% (413/1258)
and 11.1\% (214/1258) in the injury group, versus 32.8\% (413/1258) and
17.0\% (214/1258) in the adult group and 57.8\% (310/536) and 22.6\%
(121/536) in the aged group, respectively. Monocle~2 pseudotime rooted at
qNSC1/qNSC2 then progressed through pNSCs to aNSCs or neuroblasts, with
HOPX and PAX6 at the root, VIM, CD44, TNC and CHI3L1 in the middle where
pNSCs reside, SOX2, CKAP5 and RANGAP1 in aNSCs, and STMN2 in the
neuroblast terminus. Four pseudotime gene clusters captured Wnt-signalling
and regeneration at the root (TNC, SOX2, LPAR1, CLU), neurogenesis in the
middle (PTPRG, NRG3, LGI1, SLC4A4), and neuronal (MAP2, SEMA5A, SYT1,
SYN2, SYN3, MYT1L) and glial (LAMP2, PLP1, MBP, MOG) fates at the two
branch endpoints, consistent with reactivated NSCs possessing latent
potential to generate both neurons and oligodendrocytes.

To identify the molecular drivers of reactivation, we compared
neurogenic-lineage gene expression between aged and injured hippocampi.
Genes significantly up-regulated in injury covered regeneration (TNC,
GPC6, CRYAB, GBP2), autophagy/proliferation (VMP1, CHI3L1, SPP1, VIM,
ITGB1, CYR61, CD63, ACTN4, ELL2, SPOCD1), inflammation and metabolism
(IFI44L, RASGEF1B, STAT1), and metabolism alone (ACTN2, GBE1, NAMPT). Many
of these genes have rodent precedents: Tnc, Gpc6, Cryab and Gbp2 promote
neuron regeneration and synapse formation after stroke; Vmp1, Chi3l1, Spp1,
Vim and Itgb1 are linked to autophagy, proliferation and regeneration; and
Cyr61, CD63, Actn4, Ell2 and Spocd1 support proliferation, consistent with
the dormant-NSC activation and primed-quiescent-state transition described
in the mouse SVZ \cite{llorensbobadilla2015single}. The neurogenic signal
in the injured hippocampus, together with the pseudotime bifurcation into
neuronal and glial fates, is what one would predict if injury re-entered
an evolutionarily conserved reactivation program.

\subsection{Summary table of neurogenic-lineage composition}

Table~\ref{tab:lineage} summarizes the within-lineage NSC-subtype fractions
across the contexts reported in the experimental log.

\begin{table}[htbp]
\centering
\caption{Neurogenic-lineage subpopulation fractions across contexts (values
from experimental log, Figure 6D). In adult and aging groups, pNSC and
aNSC were reported as ``minimal presence'' and are therefore not quoted
as numerical fractions here.}
\label{tab:lineage}
\begin{tabular}{lcccc}
\toprule
Group & qNSC1 & qNSC2 & pNSC & aNSC \\
\midrule
Adult (31--32y) & 32.8\% (413/1258) & 17.0\% (214/1258) & minimal & minimal \\
Aging (50--68y) & 57.8\% (310/536) & 22.6\% (121/536) & minimal & minimal \\
Stroke injury (48y) & 23.4\% (413/1258) & 11.1\% (214/1258) & 17.3\% (322/1861) & 11.7\% (218/1861) \\
\bottomrule
\end{tabular}
\end{table}

\begin{table}[htbp]
\centering
\caption{Summary of newly identified human hippocampal NSC/NB markers and
the methods that surfaced them.}
\label{tab:markers}
\begin{tabular}{lll}
\toprule
Population & Markers & Method \\
\midrule
qNSC1 / qNSC2 & LRRC3B, RHOJ, SLC4A4, GLI3 & scHPF + Seurat FindAllMarkers \\
pNSC & CHI3L1, EGFR & scHPF + Seurat FindAllMarkers \\
Neuroblast (raw) & NRGN, NRN1, HECW1 & scHPF + Seurat FindAllMarkers \\
Neuroblast (filtered) & CALM3, TTC9B, NRGN, FXYD7, NRN1, GNG3, & scHPF + FindAllMarkers, \\
                      & TCEAL5, TMSB10, NEUROD2 & filtered vs.\ interneurons \\
\bottomrule
\end{tabular}
\end{table}

% Inline the figures produced by the plot agent.

\begin{figure}[htbp]
\centering
\includegraphics[width=0.95\textwidth]{figures/fig_donor_cohort.pdf}
\caption{Donor cohort and experimental design. Ten postmortem human
hippocampi spanning neonatal (postnatal day 4, $n=1$), adult (31 and 32
years, $n=2$), aging (50, 56, 60, 64 (x2), 68 years, $n=6$), and
stroke-induced injury (48 years, $n=1$) donors were dissected at the
anterior-mid hippocampus containing the dentate gyrus, then processed by
10x Genomics single-nucleus RNA sequencing in parallel with paired
immunostaining; 99{,}635 nuclei sequenced and 92{,}966 retained after
QC, yielding 16 major cell populations.}
\label{fig:cohort}
\end{figure}

\begin{figure}[htbp]
\centering
\includegraphics[width=0.95\textwidth]{figures/fig_cell_populations.pdf}
\caption{Nucleus counts per annotated cell population after quality control
(total 92{,}966 nuclei). Granule cells (GC, 24{,}671) dominate, followed by
oligodendrocytes (OLG, 15{,}796), microglia (MG, 11{,}823), astrocyte2/qNSC
(AS2/qNSC, 11{,}071), and interneurons (IN, 8{,}601); the neurogenic
progenitor fractions pNSC (2{,}798), aNSC (2{,}140), and neuroblast (NB,
2{,}607) are comparatively rare but resolvable.}
\label{fig:populations}
\end{figure}

\begin{figure}[htbp]
\centering
\includegraphics[width=0.82\textwidth]{figures/fig_lineage_fractions.pdf}
\caption{Stroke injury rebalances the neurogenic lineage toward primed and
active NSC states. Within-lineage percentages of qNSC1, qNSC2, pNSC, and
aNSC in adult, aging, and stroke-injured hippocampus. qNSC1 rises from
32.8\% in adult to 57.8\% in aging and then falls to 23.4\% after stroke,
while pNSC and aNSC --- reported as ``minimal'' in adult and aging groups
(shown as zero) --- recover to 17.3\% and 11.7\% respectively in the
injured hippocampus.}
\label{fig:lineage}
\end{figure}

\begin{figure}[htbp]
\centering
\includegraphics[width=0.92\textwidth]{figures/fig_injury_counts.pdf}
\caption{Absolute numerator / denominator counts underlying the
neurogenic-lineage fractions. Stroke-injured counts:
qNSC1 413/1258 (23.4\%), qNSC2 214/1258 (11.1\%), pNSC 322/1861 (17.3\%),
aNSC 218/1861 (11.7\%). Aged counts: qNSC1 310/536 (57.8\%), qNSC2 121/536
(22.6\%). All values taken verbatim from the experimental log.}
\label{fig:injury_counts}
\end{figure}

\section{Discussion}

Our results provide a single molecular framework in which several
previously contradictory observations about human adult hippocampal
neurogenesis become compatible. By covering neonatal, adult, aging, and
stroke-injured tissue within one atlas, we show that the neurogenic
lineage is not a monolithic on/off state but a dynamic composition in
which deep and shallow quiescence, priming, and activation coexist,
shift with age, and can be redistributed by injury. The observation
that deep-quiescent qNSC1 rises from 32.8\% (413/1258) of the neurogenic
lineage in adults to 57.8\% (310/536) in aging and then falls to 23.4\%
(413/1258) after stroke, together with the parallel recovery of pNSCs and
aNSCs from ``minimal'' in adult and aging donors to 17.3\% (322/1861)
and 11.7\% (218/1861) in the stroke donor, is the most direct quantitative
reconciliation we can offer: the reservoir persists, but its activation
status depends on life stage and pathological state.

Placing our findings against prior immunohistochemical work helps
disentangle the long-running adult-neurogenesis debate. The
Sorrells-Cipriani view that DCX-positive immature neurons are essentially
absent in adult human DG \cite{sorrells2018human,cipriani2018hippocampal}
is consistent with our finding that pNSCs and aNSCs are present only
minimally in adult and aging donors. The Boldrini-Moreno-Jim\'{e}nez-Tobin
view that hundreds or thousands of immature neurons persist into late
life \cite{boldrini2018human,morenojimenez2019adult,tobin2019human} is
consistent with our observation that a subset of neuroblasts is retained
into adulthood (longer than in rodents), with our detection of reactivated
NSCs after injury, and with the finding by Terreros-Roncal and colleagues
that neurodegenerative disease substantially alters the neurogenic niche
\cite{terrerosroncal2021impact}. The two positions need not be mutually
exclusive: if healthy adult tissue carries primarily quiescent cells that
occasionally transition to shallow quiescence, while disease and injury
drive the pool toward the primed/active state, both camps are partially
correct. The discrepancy would then come down to sample composition ---
specifically whether donor brains carry comorbidities or pre-mortem
insults that push the pool toward reactivation.

The interneuron-contamination problem we demonstrate has direct
methodological implications. DCX, PROX1, STMN2, CALB2, and several other
canonical neuroblast markers are widely expressed in human GABAergic
interneurons, consistent with an earlier observation in human hippocampus
\cite{franjic2022transcriptomic}, and co-immunostaining confirmed that
SST-positive interneurons in primate hippocampus can be DCX-positive.
Any study that identifies ``new neurons'' in adult human hippocampus on
the basis of these markers alone therefore risks scoring mature
interneurons. Our combined scHPF-plus-FindAllMarkers pipeline followed by
explicit interneuron-expression filtering yielded human neuroblast markers
(CALM3, TTC9B, NRGN, FXYD7, NRN1, GNG3, TCEAL5, TMSB10, NEUROD2) that
were validated across adult and aging samples, and we believe this class
of filters should become a standard component of adult-neurogenesis
pipelines going forward. The complementary set of markers for qNSCs
(LRRC3B, RHOJ, SLC4A4, GLI3) and pNSCs (CHI3L1, EGFR) likewise offers a
finer-grained alternative to the mostly astrocytic profile (GFAP, ALDH1L1,
VIM) that currently dominates human NSC identification.

The stroke-induced reactivation signature we observe aligns with the
dormant-NSC activation program described in the mouse subventricular zone
after ischemia \cite{llorensbobadilla2015single} and with the broader
framework of injury-modulated NSC quiescence in the adult rodent brain
\cite{obernier2019neural,urban2019quiescence}. The up-regulation of
regeneration (TNC, GPC6, CRYAB, GBP2), autophagy and proliferation
(VMP1, CHI3L1, SPP1, VIM, ITGB1, CYR61, CD63, ACTN4, ELL2, SPOCD1),
inflammation/metabolism (IFI44L, RASGEF1B, STAT1), and metabolism (ACTN2,
GBE1, NAMPT) genes in the injured neurogenic lineage mirrors these
rodent precedents. At the same time, the near-absence of neuroblasts in
the injury sample --- despite the strong reactivation of pNSCs and aNSCs
--- likely reflects the two-day post-stroke interval, too short for
injury-activated NSCs to mature into detectable neuroblasts in a species
whose neuroblast-maturation window exceeds six months in macaques
\cite{kohler2011maturation}. A complete reactivation-to-maturation
transition may simply not have time to occur in acute stroke donors.

Several limitations must be acknowledged. The neonatal and stroke cohorts
contain a single donor each ($n=1$), and lifestyle and health parameters
for all donors are incompletely documented, so the quantitative
percentages should be read as representative rather than as
population-level estimates. snRNA-seq cannot fully separate cell types
with very similar transcriptomes --- qNSCs and astrocytes being the most
salient example --- and therefore our gene-set-score separation of those
populations, while corroborated by marker genes and cell-cycle state,
is not definitive. The very rare NSCs captured in aged and injured
samples also lack protein-level functional validation beyond the
immunostainings we performed. Future progress will likely come from
integrating snRNA-seq with epigenomic, proteomic, and metabolomic
profiling at single-cell resolution \cite{ayhan2021resolving,habib2017massively},
from non-invasive human lineage-tracing technologies, and from expanding
the donor cohort to multiple neonatal and stroke samples so that the
$n=1$ observations reported here can be replicated. Together, the atlas
and analytical pipeline we describe should serve as a reference point
for such extensions, and the molecular markers we identify should make
the identification and quantification of reactivated NSCs more reliable
in future human hippocampal studies.

\bibliographystyle{unsrtnat}
\bibliography{references}

\end{document}
